%% SCRAP: archive/research/research-outline
%% SOURCE: docs/working/archive/research/research-outline.md
%% STATUS: WORKING
%% FITS: vol3-research/ if relevant
%% EDITORIAL: lifted — prose rewritten to press voice

\section{Research Paper Outline: Physics-Driven VM Optimisation}
\label{sec:research-outline}

\subsection{Working Title}

\textit{Physics-Inspired Threshold-Based Runtime Optimisation in Virtual Machines}

Alternative titles under consideration:
\begin{itemize}
  \item ``Eliminating JIT: Automatic VM Optimisation Without Code Generation''
  \item ``Real-Time Metrics-Driven Optimisation in Microkernel-Compatible VMs''
  \item ``Bayesian Inference at Nanosecond Scale: Pure Fixed-Point Optimisation Decisions''
\end{itemize}

\subsection{Target Format and Scope}

ACM SIGPLAN format; estimated 25--34 pages plus figures, tables, references,
and optional appendix.

\subsection{Section Structure}

\subsubsection{1. Introduction (4 pages)}

Opens with the observation that VM optimisation traditionally requires
offline profiling, JIT compilation, or expert manual tuning, then presents
the alternative: automatic, real-time, metrics-driven optimisation requiring
neither code generation nor floating-point arithmetic.

Four key questions are posed:
\begin{enumerate}
  \item How can meaningful VM performance improvement be achieved without JIT?
  \item Can optimisation decisions be made automatically in real time?
  \item Can a pure-integer statistical framework validate optimisation impact?
  \item Can such a system remain compatible with formal verification?
\end{enumerate}

Five primary contributions are claimed: the physics-inspired framework,
threshold-based automatic decision logic, real-time Bayesian inference in
Q48.16 fixed-point, a proven 1.78$\times$ speedup with statistical
validation, and nine additional optimisation opportunities using identical
infrastructure.

\subsubsection{2. Motivation and Problem Analysis (3 pages)}

Surveys the three classical approaches: offline profiling (simple but
reactive), JIT compilation (high performance but complex, hard to verify,
and microkernel-incompatible), and manual tuning (verifiable but not
scalable). Poses the research question: can a 1.5--2.0$\times$ improvement
be achieved without any of those approaches while remaining verifiable and
L4Re-compatible?

\subsubsection{3. Related Work (4 pages)}

Covers adaptive optimisation, metrics-driven systems, physics-inspired
computing, fixed-point arithmetic and formal verification, FORTH and
stack-based VMs, and microkernel-compatible systems. See
Section~\ref{sec:literature-review} for full detail.

\subsubsection{4. System Design and Architecture (5 pages)}

Describes the four layers: metrics collection (\texttt{execution\_heat},
\texttt{temperature\_q8}, nanosecond-precision latency in Q48.16), decision
logic (threshold check: if \texttt{execution\_heat} $> 50$ then promote),
the 32-entry LRU hot-words cache, and the Bayesian inference engine
(Beta-Binomial posterior, 95\% and 99\% credible intervals, all arithmetic
in Q48.16).

\subsubsection{5. Experimental Methodology (4 pages)}

Platform: x86\_64, 100\% assembly optimisations plus LTO. VM: StarForth
FORTH-79 in strict ANSI~C99. Benchmark: dictionary lookup with 23 common
FORTH words, 100,000 samples, \texttt{CLOCK\_MONOTONIC\_RAW}. Two
build variants compared: \texttt{ENABLE\_HOTWORDS\_CACHE=1} versus
\texttt{ENABLE\_HOTWORDS\_CACHE=0}. Three independent runs for
reproducibility validation.

\subsubsection{6. Results (5 pages)}

Core result: 1.78$\times$ speedup (bucket mean 56.237\,ns vs cache mean
31.543\,ns). 95\% credible interval [1.75$\times$, 1.81$\times$]. Cache
hit rate 35.64\%. Ten words automatically promoted; zero manual tuning.
Run-to-run standard deviation 0.0015$\times$. Memory overhead less than
1\,KB; per-lookup overhead less than 2 CPU cycles.

\subsubsection{7. Analysis and Discussion (4 pages)}

Explains why the approach succeeds: execution frequency is predictable
(Zipfian distribution), threshold logic is robust, and Q48.16 arithmetic
is sufficient for nanosecond-precision timing without floating-point error
accumulation. Compares against JIT. Identifies nine additional optimisation
opportunities (stack fusion, vocabulary reordering, return-stack prediction,
and six others) with a projected cumulative improvement of 5--8$\times$.

\subsubsection{8. Limitations and Future Work (3 pages)}

Current limitations: restricted to dictionary lookup, single-threaded VM,
x86\_64 only, 100K sample workload. Future work: extended optimisations
(Phase~2), adaptive thresholds (Phase~3), Isabelle/HOL proof of cache
correctness (Phase~4), and ARM64 and L4Re validation (Phase~5).

\subsubsection{9. Conclusion (2 pages)}

Summarises five contributions and their broader impact: formal verification
is compatible with performance; JIT is not necessary for measurable speedup;
physics-inspired metrics represent a new research direction; the approach is
practical for microkernel ecosystems.

\subsection{Key Figures}

\begin{enumerate}
  \item System architecture: metrics $\to$ decision $\to$ cache.
  \item Latency comparison histogram (cache vs bucket).
  \item Speedup with credible intervals.
  \item Cache contents and \texttt{execution\_heat} distribution.
  \item Word promotion timeline (frequency vs time).
\end{enumerate}

\subsection{Key Tables}

\begin{enumerate}
  \item Lookup performance summary (cache vs bucket latency).
  \item Bayesian posterior distributions.
  \item Comparison with JIT and offline profiling approaches.
\end{enumerate}
